\documentclass{article}
\usepackage[spanish]{babel}
\usepackage{tikz}
\usetikzlibrary{babel} % REGLA 1: Soluciona conflictos de flechas con Babel en español
\usepackage{circuitikz}
\usepackage{siunitx}   % REGLA 6: Formato estándar para números y unidades
\usepackage{listings}  % REGLA 4: Inclusión segura de código fuente
\usepackage{xcolor}

\lstset{
    basicstyle=\ttfamily\small,
    language=Python,
    frame=single,
    backgroundcolor=\color{gray!5}
}

\begin{document}

\section{Prueba de Memoria Evolutiva (LaTeX)}

% REGLA 5: Escape de caracteres especiales
Este documento verifica que los caracteres especiales como el guion\_bajo y los porcentajes (100\%) no rompan la compilación y funcionen correctamente en el texto plano.

\subsection{Diagrama TikZ / Circuitikz}
% REGLA 6: Uso de siunitx (\qty) en el texto
El siguiente circuito cuenta con una resistencia de \qty{10}{\kilo\ohm} y un capacitor de \qty{100}{\nano\farad}.

\begin{center}
\begin{circuitikz}
    % REGLA 7: Etiquetas de componentes estrictamente en modo matemático ($...$)
    \draw (0,0) to[R, l={$R_1 = \qty{10}{\kilo\ohm}$}] (2,0)
          to[C, l={$C_1 = \qty{100}{\nano\farad}$}] (4,0);
    
    % REGLA 2: Parseo seguro de coordenadas separando explícitamente \x/\y
    \foreach \x/\y in {0/-1.5, 2/-1.5, 4/-1.5} {
        \node[circle, fill=blue, inner sep=1.5pt] at (\x,\y) {};
    }
    
    % Flecha con TikZ que fallaría sin \usetikzlibrary{babel}
    \draw[->, thick, red] (1, -0.5) -- (1, -1.2) node[midway, right] {Nodo test};
\end{circuitikz}
\end{center}

\end{document}